% problem2.tex
\section*{Problem 2 — Horner's Method}

There is a more efficient algorithm for evaluating polynomials than the conventional method from Problem 1. Horner's method evaluates
\[
a_n x^{n} + a_{n-1}x^{n-1} + \cdots + a_1 x + a_0
\]
at $x=c$ using the pseudocode:
\[
\begin{aligned}
&\textbf{procedure } \mathrm{HORNER}(c, a_0, a_1, \ldots, a_n : \text{ real numbers})\\
&\qquad y \leftarrow a_n\\
&\qquad \textbf{for } i = 1 \text{ to } n \textbf{ do}\\
&\qquad\qquad y \leftarrow y \cdot c + a_{n-i}\\
&\qquad \textbf{end for}\\
&\qquad \textbf{return } y\\
&\textbf{end procedure}
\end{aligned}
\]

\subsection*{(a) Evaluate $3x^2 + x + 1$ at $x=2$ by tracing each step}
Here $n=2$, $c=2$, and $(a_0,a_1,a_2)=(1,1,3)$.

\begin{center}
\begin{tabular}{c|l|c}
\textbf{Step} & \textbf{Action} & \textbf{$y$} \\
\hline
1 & $y \leftarrow a_2 = 3$ & $3$ \\
\hline
2 & $i=1$: $y \leftarrow y\cdot c + a_{2-1} = 3\cdot 2 + 1$ & $7$ \\
3 & $i=2$: $y \leftarrow y\cdot c + a_{2-2} = 7\cdot 2 + 1$ & $15$ \\
\end{tabular}
\end{center}

\noindent\textbf{Return: } $y=15$. Thus $3(2)^2 + 2 + 1 = 15$.

\subsection*{(b) Exact operation counts}
Each loop iteration performs:
\begin{itemize}
  \item one multiplication: $y \cdot c$,
  \item one addition: $+\, a_{n-i}$.
\end{itemize}
With $n$ iterations total, the exact counts are:
\[
\boxed{\text{Multiplications} = n}
\qquad
\boxed{\text{Additions} = n}.
\]